% !TEX root = ../main_ver2_bench_only.tex
% NOTE: \section*{Limitations} is emitted by main.tex; this file starts directly with content.
\label{sec:limitations}

\noindent
\textbf{Synthetic diagnostic tasks.}\quad
Step~0 evaluates each skill using a synthetic task generated by a language model rather than drawn from a fixed clinical task distribution.
While this design provides controlled, per-skill isolation, it may not fully represent the diversity and complexity of real clinical workflows.
We mitigate this concern through our benchmark integration analysis (\Cref{sec:experiments}), which validates skill relevance against authentic questions from 12 established medical QA benchmarks.

\textbf{Binary 8-dimensional scoring.}\quad
Our evaluation framework scores each execution trace on eight binary dimensions.
This design supports clear pass/fail attribution and aggregation at scale, but it sacrifices continuous granularity---for instance, a partially correct parameter schema receives the same score as a completely wrong one.
Future work could augment binary verdicts with ordinal or continuous rubrics for finer-grained diagnostics.

\textbf{LLM-as-judge bias.}\quad
The scoring pipeline relies on a language model as the judge.
We mitigate potential bias through role separation (the judge model differs from the executor), rule-based signal extraction, and human spot-checks on a stratified sample.
Nevertheless, systematic expert review---particularly by domain specialists in medicine---remains necessary before drawing clinical conclusions from the evaluation scores.

\textbf{Single-model primary evaluation.}\quad
Step~0 results reported in this paper are primarily obtained with GPT-5.1.
While backup runs on Claude Sonnet~4.5 confirm broad trend consistency, comprehensive multi-model evaluation---including open-weight models---is planned as future work.
